\documentclass[italian,12pt]{article} %tipo di documento

\def\Title{Verbale interno}
\def\Subtitle{Riunione interna settimanale}
\def\Author{7Last}
\def\Date{2024-05-22}
\def\Version{v1.0}

%pacchetti extra da scaricare dblfloatfix, fancyhdr
\usepackage[left=2cm, right=2cm, bottom=3cm, top=3cm]{geometry}
\usepackage{fancyhdr}%creazione header-footer
\usepackage{graphicx} %serve per inserire immagini
\graphicspath{ {../../../logo/} }
%\usepackage{dblfloatfix} %serve per posizionare gli elementi dove si vuole
\usepackage[hidelinks]{hyperref} %serve per i link
\usepackage{tikz}
\usepackage{tgadventor} % font
\usepackage[useregional=numeric,showseconds=true,showzone=false]{datetime2}
\usepackage{caption}

\usepackage{hyperref}
\usepackage{tocloft}
\usepackage{titlesec}
\usepackage{color}
\usepackage{ulem}
\usepackage{pgfplots}
\usepackage{pgf-pie}
\usepackage[italian]{babel}
\usepackage{comment}
\usepackage{tabularx}
\usepackage{longtable}
\usepackage{float}
\usepackage{amsmath}

\linespread{1.2}
\captionsetup[table]{labelformat=empty}

\renewcommand{\contentsname}{Indice}%imposto il nome dell'indice
\renewcommand\familydefault{\sfdefault}
%-------------------INIZIO DOCUMENTO--------------
\begin{document}

\newgeometry{left=2cm,right=2cm,bottom=2.1cm,top=2.1cm}
\begin{titlepage}
    \vspace*{.5cm}

    \vspace{2cm}
    {
        \centering
        {\bfseries\huge \Title\par}
        \bigbreak
        {\bfseries\large \Author\par}
        \bigbreak
        {\Version\par}
        \vfill

        \begin{tikzpicture}[remember picture,overlay]

            \fill[blue!20!black] (current page.south west) -- ++(10cm,0) -- ++(-10cm,15cm) -- cycle;
            \fill[orange] (current page.south east) -- ++(-18cm,0) -- ++(21.6cm,22cm) -- cycle;

            \clip (0,-3cm) circle (2.5cm) node (current page.center) {\includegraphics[width=5cm]{logo.jpg}};
        \end{tikzpicture}

    }

    \vfill

\end{titlepage}

\restoregeometry






















\newpage

%---------------header------------------%
\pagestyle{fancy}
\fancyhead{} % pulizia degli header
\lhead{%
\begin{tikzpicture}
    \clip (0,0) circle (0.5cm);
    \node at (0,0) {\includegraphics[width=1cm]{logo}};
\end{tikzpicture}%
}
\chead{\vspace{\fill}\Title\vspace{\fill}}
\rhead{\vspace{\fill}\Version\vspace{\fill}}
%quad è una spaziatura
%---------------------------------------%



%-----------tabella versioni-----------%
\begin{table}[!h]
	\caption{Versioni}
	\begin{center}
		\begin{tabular}{ c c c p{9cm}}
			\hline \\[-2ex]
			Ver. & Data & Autore & Descrizione \\
			\\[-2ex] \hline \\[-1.5ex]
			1.0 & 18/03/2024 & Elena Ferro &  Approvazione documento\\
			0.3 & 18/03/2024 & Matteo Tiozzo & Prima revisione \\
			0.2 & 15/03/2024 & Valerio Occhinegro& Sistemazione ortografica e di linguaggio  \\
			0.1 & 08/03/2024 & Leonardo Baldo& Prima stesura del documento \\
			\\[-1.5ex] \hline
		\end{tabular}
	\end{center}
\end{table}
%---------------------------------------%
\newpage

\tableofcontents

\newpage
\section{Partecipanti}

\textbf{Sede della riunione}: piattaforma meet\\
\textbf{Orario di inizio}: 15:00\\
\textbf{Orario di fine}: 15:30\\


\begin{flushleft}
	\begin{table}[!h]
	\begin{tabular}{ |c|c|c| } 
		\hline
		\textbf{Componente} & \textbf{Ruolo} & \textbf{Presenza} \\
		\hline 
		Leonardo Baldo 		& Redattore & Presente \\ 
		Antonio Benetazzo 	& Responsabile & Presente \\
		Elena Ferro 		& Verificatore & Presente \\
		Valerio Occhinegro 	& Redattore & Presente \\
		Raul Seganfreddo 	& Amministratore & Presente \\
		Matteo Tiozzo 		& Verificatore & Presente \\ 
		Davide Malgarise 	& Redattore & Assente \\
		\hline
	\end{tabular}

	\vspace{1cm}

	\begin{tabular}{ |c|c| } 
		\hline
		\textbf{Nome} & \textbf{Ruolo} \\
		\hline 
		Fabio Pallaro & Rappresentante dell'azienda \\
		Daniele Zorzi & Rappresentante dell'azienda \\
		Andrea Dorigo & Rappresentante dell'azienda \\

		\hline
	\end{tabular}
\end{table}

\end{flushleft}


\newpage

\begin{flushleft}
\section{Verbale dell'incontro}
\subsection{Linguaggi di programmazione e tecnologie utilizzate}
	\textbf{Domanda}: è possibile utilizzare linguaggi di programmazione differenti da Python per la simulazione dei dispositivi IOT?\\
	È possibile utilizzare altri linguaggi di programmazione, l'utilizzo di Python è consigliato per la facilità del codice e per la popolarità stessa del linguaggio.


\subsection{Quantità di dati da gestire}
	\textbf{Domanda}: la quantità dei dati da gestire sarà paragonabile a quella effettivamente prodotta dai sensori?\\
	L'archittettura dovrà essere scalabile, ma per questo progetto si tratterà di una simulazione, sarà dunque sufficiente analizzarne un numero simbolico.

\subsection{Elementi di complessità}
	\textbf{Domanda}: quali sono gli elementi di complessità maggiore del progetto?\\
	Gli elementi che richiederanno maggiore impegno sono l'integrazione delle varie tecnologie tra loro, la creazione dell'architettura (networking, containerizzazione) e la configurazione (devops) dei prodotti. La parte di programmazione sarà limitata allo script python per la generazione dei dati.

\subsection{Strumenti consigliati}
	\textbf{Domanda}: quali strumenti consigliate per la gestione del progetto?\\
	Per la gestione del progetto l'azienda consiglia l'utilizzo di ClickUp, Discord e Trello.

\subsection{Disponibilità e organizzazione Sync Lab S.r.l.}
	\textbf{Domanda}: quale supporto verrà fornito dall'azienda? Come saranno organizzate le riunioni e gli incontri?\\
	L'azienda fornirà una formazione per le tecnologie che andranno utilizzate e metterà a disposizione un tutor per la risoluzione di eventuali problemi. Sottolinea che sarà disponibile ad effettuare un incontro di routine una volta a settimana, da concordare con il gruppo, e sarà disponibile in caso di necessità. 

\end{flushleft}

\begin{table}[b]
	\begin{tabular}{@{}p{.5in}p{4in}@{}}
		Data:  & \hrulefill \\
			   &     		\\
			   &     		\\
		Firma: & \hrulefill \\
	\end{tabular}
	\end{table}

\end{document}

























